\begin{abstract}
A short superpermutation must make extensive use of large overlaps between
successive permutations. These overlaps organize the permutations into
runs, and earlier lower bounds control how efficiently such runs can be
joined. We strengthen this part of the argument. The patterns that attain
the earlier estimates cannot keep repeating: their action on the symbols
eventually brings them back to a rotation class already visited. In one
case, two repetitions already have an entire run in common. Counting the
necessary departures from these patterns gives an explicit lower bound
valid for every alphabet size at least eight. It improves the preceding
general formula at every size at least thirteen not divisible by three.
The leading correction above $K!+(K-1)!+(K-2)!$ remains asymptotic
to $2(K-4)!$. The exact formulas and their asymptotic comparisons have
formal proofs in Lean.
\end{abstract}

\noindent\textbf{AI disclosure.} The research presented here, the Lean
formalization, and this document were generated entirely by AI systems
under Jay Pantone's direction. The inherited mathematical work and
software retain their original credits. Independent AI review is not
human peer review.

\section{Introduction}

A \emph{superpermutation} on $K$ symbols is a word containing each of the
$K!$ permutations as a consecutive subword. For example, $123121321$
contains all six permutations on three symbols. Write $S(K)$ for the
length of a shortest superpermutation and put
\[
 F_3(K)=K!+(K-1)!+(K-2)!.
\]
The lower bounds studied here exceed $F_3(K)$ by a quantity asymptotic
to $2(K-4)!$. We improve the smaller terms by examining when the local
counting estimates can be tight. The resulting inequalities hold at every
alphabet size in their stated range.

The difficulty is that a pattern may be efficient for a few steps without
remaining efficient for very long. One can see the simplest instance
directly in a word. Appending the first symbol of a permutation produces
its next cyclic rotation, at a cost of only one letter. Repeating this
operation is inexpensive, but after $K$ steps it returns to the original
permutation. A covering word must eventually do something else.

The same phenomenon occurs further into the lower-bound argument.
After grouping rotations together, and then grouping those classes into
runs, the relevant objects are chains of almost full runs. The earlier
proof bounds the length of each chain and restricts which lengths may
follow one another. Its best average can nevertheless be attained by a
repeating sequence of chain lengths. Such a sequence does not determine
fresh permutations indefinitely. The corresponding relabeling of the
symbols has finite order, and when it returns, so does a previously
visited class.

This gives a way to improve the average. Each interruption of the
repeating pattern loses some of the length allowed by the old estimate.
If there is a uniform bound on the distance between interruptions, these
losses accumulate. There are two parts to making this observation useful:
finding a sufficiently short return, and proving that all other patterns
really do lose enough to pay for the interruptions. Most of the work in
this paper concerns the second part.

Our explicit bound is stated in Theorem~\ref{thm:main}. It applies for
every $K\ge8$, and incorporates the earlier refinement for $K$ divisible
by three. In each of the other two residue classes it strictly improves
the preceding general formula for every $K\ge13$. At $K=13$ and
$K=14$, for example, the increases are $12$ and $58$, respectively. The improvements hold uniformly
in the alphabet size; no separate search is used to prove an instance.

We use the geometric and forest lemmas of the lower-bound framework
retained in the accompanying proof library~\cite{pantone-framework}.
These lemmas reduce an arbitrary word to auxiliary paths of complete
rotation classes. The new estimates in this paper concern those paths.
The next section explains why counting deficits in such paths gives a
word-length bound. We then describe the chains that make the count
difficult. A fourteen-symbol example introduces the return argument
before we prove it in general. The explicit numerical formula is given
after these arguments, when its parameters have a combinatorial meaning.
The coordinate calculations and the endpoint count for the forest
reduction are given in the appendices.
